\section{La experiencia de Li Auto: datos, presión organizacional y el proyecto Weicheng}

Las secciones anteriores trataron las rutas técnicas; esta sección vuelve a la práctica de Li Auto (理想汽车). Cuando Lang Xianpeng (郎咸朋) se incorporó a Li Auto en 2018, ya había discutido con Li Xiang (李想) que lo más importante en la conducción autónoma son los datos: el talento se puede reclutar y la capacidad de cómputo se puede comprar, pero los datos reales de la flota, los casos de cola larga y los escenarios de los usuarios no se pueden comprar fuera. La ventaja de un fabricante de automóviles de producción en serie está en tener flota, usuarios y retroalimentación de caminos reales; la dificultad es que debe iterar rápido bajo presiones de seguridad, costo, experiencia y organización.

\begin{figure}[H]
\centering
\includegraphics[width=0.96\textwidth]{ideal-av-data-loop.png}
\caption{Ciclo cerrado de datos de la conducción inteligente de Li Auto: la flota de producción en serie, los casos de cola larga, el entrenamiento y la retroalimentación del despliegue forman un volante de inercia. Diagrama conceptual propio, elaborado a partir de la conversación de 00:41:14--00:50:50 y 01:26:40--01:35:51.}
\end{figure}

\begin{knowledgebox}{Lectura de la figura: la ventaja de datos del fabricante proviene de un ciclo cerrado real}
De Fleet, Cases, Label/Train y Deploy a Feedback. La flota de producción en serie no es un almacén de datos, sino un ciclo cerrado que descubre continuamente problemas de cola larga, entrena modelos, los despliega por OTA y vuelve a recoger retroalimentación. El techo de la conducción autónoma lo determinan conjuntamente la calidad de los datos y la capacidad de cómputo para entrenamiento.
\end{knowledgebox}

\subsection{Privacidad de datos y datos del exterior del vehículo}

Esta sección agrega un límite importante. En la entrevista, Lang Xianpeng enfatizó que Li Auto recoge datos del exterior del vehículo y no recoge información biométrica de los usuarios, como rostros o voces dentro del vehículo; además, al transmitir los datos del exterior se procesan los rostros y las placas. Esto no es un detalle de relaciones públicas, sino la base para que el ciclo cerrado de datos de la conducción autónoma pueda operar a largo plazo: sin cumplimiento normativo y sin la confianza de los usuarios, el volante de inercia de datos reales no puede sostenerse.

\begin{importantbox}{Experiencia práctica: el ciclo cerrado de datos requiere confianza como condición previa}
La conducción autónoma necesita una gran cantidad de datos reales de caminos, pero los datos reales no son un recurso sin límites. La captura en el exterior del vehículo, la anonimización, el control de permisos, la limitación de usos y el almacenamiento seguro determinan si los datos de la flota pueden convertirse en un activo de largo plazo.
\end{importantbox}

\subsection{De Baidu a Li Auto: cómo la formación en una ruta cambia la organización}

Esta sección complementa la trayectoria personal de Lang Xianpeng. En 2013 participó en Street View y en mapas de alta precisión dentro del equipo de Baidu Maps, y después en la colaboración de conducción autónoma con BMW; en esa etapa se formaron capacidades en mapas, topografía, percepción y sistemas de vehículos de ingeniería. Al incorporarse a Li Auto en 2018 se enfrentó a un fabricante de producción en serie: los vehículos se venden a usuarios reales, las funciones deben llegar por OTA a una gran cantidad de vehículos y los problemas regresan de los caminos reales en forma de casos de cola larga. La diferencia entre estas dos etapas explica por qué insiste una y otra vez en los datos y en el ciclo cerrado de la producción en serie.

\begin{knowledgebox}{Nota de clase: el vehículo de investigación, el Robotaxi y el vehículo de producción en serie son tres organizaciones distintas}
El vehículo de investigación busca la validación técnica; el Robotaxi busca poder operar dentro de una zona delimitada; el vehículo de producción en serie busca costo, estabilidad, experiencia de usuario y mantenimiento durante todo su ciclo de vida. Las rutas técnicas son parecidas, pero los objetivos organizacionales son completamente distintos.
\end{knowledgebox}

\noindent\begin{tabularx}{\textwidth}{p{0.18\textwidth}p{0.36\textwidth}X}
\toprule
Etapa & Activos clave & Formación para la conducción autónoma \\
\midrule
Baidu Maps/Street View & Vehículos de captura, topografía, anonimización, mapas de alta precisión & Comprender los datos de caminos y la infraestructura de mapas. \\
Baidu ADU/colaboración con BMW & Vehículos de ingeniería, pruebas L3/L4, mapas de alta precisión & Conocer la ruta determinista temprana y las limitaciones de los vehículos de ingeniería. \\
Vehículos de producción en serie de Li Auto & Flota de usuarios, OTA, datos del exterior del vehículo, recursos organizacionales & Construir un ciclo cerrado de datos reales y la conciencia de responsabilidad de la producción en serie. \\
\bottomrule
\end{tabularx}

\subsection{El proyecto Weicheng y la presión organizacional}

Esta sección aborda la historia organizacional de la segunda mitad de la entrevista. El proyecto «Weicheng» (卫城) se relata como una experiencia de alta presión: la negociación con proveedores, la elección de ruta técnica, los plazos de producción en serie, la presión de Li Xiang y la dignidad del equipo se entrelazan. Cuando Lang Xianpeng habla de «morir de pie», lo que hay detrás es un equipo de tecnología de frontera que, en un momento decisivo, no quiso asegurar la seguridad de corto plazo arrodillándose ante proveedores externos, sino que eligió asumir el riesgo organizacional de cambiar de ruta.

\begin{figure}[H]
\centering
\includegraphics[width=0.86\textwidth]{weicheng-project.png}
\caption{Campo de presiones del proyecto Weicheng: la ruta técnica, los proveedores, la dignidad organizacional y los plazos de producción en serie presionan al mismo tiempo. Diagrama conceptual propio, elaborado a partir de la conversación de 01:32:23--01:35:51.}
\end{figure}

\begin{knowledgebox}{Lectura de la figura: la elección de ruta técnica se convierte en presión organizacional}
En el centro de la figura está Project Weicheng, rodeado por tecnología, proveedores, organización, producción en serie, liderazgo y equipo. Al leer esta figura hay que ver que la conducción autónoma no es una competencia de modelos de laboratorio: la elección de ruta afecta la entrega de la producción en serie, la relación con la cadena de suministro, la moral del equipo y la estrategia de la empresa.
\end{knowledgebox}

\subsection{Liderazgo de alta presión y ejecución de ingeniería}

Esta sección trata con cautela la discusión relacionada con Li Xiang. La entrevista menciona que Li Xiang se enojaba con el equipo de conducción inteligente y que no tiene muy buen carácter, y también menciona el objetivo detrás de esa presión: que el equipo completara la transición técnica y la entrega de la producción en serie dentro de una ventana de tiempo crítica. Al organizar este material como curso no se evalúa el carácter de la persona, sino el mecanismo organizacional: si la alta presión trae objetivos más claros, una asignación de recursos más rápida y una ejecución más sólida, y si al mismo tiempo evita el desgaste excesivo, la cultura del miedo y las concesiones en seguridad.

\begin{warningbox}{La alta presión no es correcta por sí misma}
Las organizaciones de tecnología de frontera necesitan objetivos firmes y una ejecución sólida, pero la conducción autónoma implica responsabilidad sobre la seguridad, y la alta presión no puede sustituir la validación rigurosa. Una buena presión organizacional debe impulsar que los problemas salgan a la luz, que los recursos se concentren y que las decisiones se tomen rápido, no obligar al equipo a saltarse los límites de seguridad.
\end{warningbox}

\subsection{Negociación con proveedores y capacidad de desarrollo propio}

Esta sección sitúa las expresiones «arrodillarse a suplicarle» y «morir de pie» en su contexto industrial. Los fabricantes de automóviles con conducción autónoma suelen depender de proveedores de chips, sensores, algoritmos, mapas, cadenas de herramientas e integración. Depender de proveedores puede acelerar la implementación temprana, pero si las capacidades centrales se subcontratan durante mucho tiempo, en los momentos decisivos el fabricante pierde la iniciativa sobre su ruta. La emoción con la que Lang Xianpeng habla del proyecto Weicheng no es solo cuestión de carácter personal, sino de un equipo que elige entre «seguir dependiendo de soluciones externas» y «arriesgarse a construir capacidad de desarrollo propio».

\begin{importantbox}{Las capacidades centrales no pueden depender solo de compras}
La barrera de largo plazo de la conducción autónoma está en el ciclo cerrado de datos, modelos, evaluación, despliegue en el vehículo y responsabilidad sobre la seguridad. Los proveedores externos pueden aportar herramientas y componentes, pero si el fabricante no domina el criterio central y la capacidad de iteración, le será muy difícil cambiar de rumbo con rapidez en la era de extremo a extremo.
\end{importantbox}

\begin{figure}[H]
\centering
\includegraphics[width=0.96\textwidth]{av-org-leadership.png}
\caption{Ciclo cerrado organizacional de la conducción autónoma: el criterio técnico, la presión de la dirección, la ejecución del equipo y la responsabilidad sobre la seguridad deben cumplirse al mismo tiempo. Diagrama conceptual propio, elaborado a partir de la conversación de 01:26:40--02:00:20.}
\end{figure}

\begin{knowledgebox}{Lectura de la figura: el ciclo cerrado organizacional es tan importante como el técnico}
El equipo de conducción autónoma necesita juzgar la ruta técnica, soportar la presión de la alta dirección, completar la ejecución de ingeniería y después construir confianza mediante la evaluación de seguridad y la experiencia de usuario. Si cualquier eslabón se rompe, a la ruta técnica le será muy difícil llegar al vehículo de producción en serie.
\end{knowledgebox}

\subsection{Resumen de la sección}

La experiencia de Li Auto muestra que la competencia entre empresas de conducción autónoma no se limita a modelos y algoritmos, sino que también abarca el cumplimiento normativo de los datos, el ciclo cerrado de la flota, la negociación en la cadena de suministro, la resiliencia organizacional y la validación de seguridad. La conducción inteligente de producción en serie se forma con el entrenamiento conjunto de un sistema técnico y un sistema organizacional.
